% Generated from project outputs. Do not edit by hand.
\begin{table}[!htbp]
\centering
\begin{threeparttable}
\caption{Including remittances by subsample: 2016-2019}
\label{tab:descriptive-2016-2019-subsamples-y4}
\small
\setlength{\tabcolsep}{4pt}
\renewcommand{\arraystretch}{1.15}
\begin{tabularx}{\linewidth}{@{}Xrrrrr@{}}
\toprule
Sample & N & Mean t & SD t & Mean t+1 & SD t+1 \\
\midrule
General & 276 & 2,866.2 & 1,954.8 & 2,907.3 & 1,870.3 \\
Men & 276 & 3,165.2 & 2,352.4 & 3,185.5 & 2,238.1 \\
Women & 276 & 2,353.0 & 1,541.1 & 2,441.5 & 1,573.3 \\
Urban & 276 & 3,143.4 & 1,995.0 & 3,165.8 & 1,890.8 \\
Rural & 271 & 2,170.9 & 1,515.3 & 2,297.9 & 1,621.2 \\
Indigenous & 271 & 2,132.5 & 1,426.9 & 2,180.2 & 1,395.6 \\
Non-Indigenous & 276 & 3,087.8 & 2,004.1 & 3,146.2 & 1,934.4 \\
Bottom 99 percent & 276 & 2,618.4 & 1,455.7 & 2,668.1 & 1,412.6 \\
Bottom 90 percent & 276 & 2,176.8 & 1,036.0 & 2,214.8 & 989.2 \\
Bottom 50 percent & 272 & 1,251.3 & 564.8 & 1,299.7 & 548.7 \\
\bottomrule
\end{tabularx}
\begin{tablenotes}[flushleft]
\footnotesize
\item Monthly income in quetzales. Unweighted descriptive statistics of survey-weighted cohort means. N counts cohort-transition rows, not individuals or independent cohorts. SD is dispersion across cohort means, not the standard error of a mean. Both incomes must be observed. Subsamples overlap; bottom-income groups are cumulative.
\end{tablenotes}
\end{threeparttable}
\end{table}
